\chapter[Degree-Five Quartic Theorem]{The Degree-Five Quartic Theorem: A Boundary of Relational Resolution}

\section{Why Degree Five Is Different}

The $\ell=5$ representation is the clearest technical example in the book because its quartic energy develops a minimum manifold rather than an isolated shape. This is not a numerical accident and not a metaphor imported from the later consciousness discussion. It follows from the structure of the invariant algebra and can be demonstrated exactly at an explicit regular point.

For $V_5$, the invariant dimensions through degree six are

\[
\dim\operatorname{Inv}_2(V_5)=1,
\qquad
\dim\operatorname{Inv}_3(V_5)=0,
\]

\[
\dim\operatorname{Inv}_4(V_5)=2,
\qquad
\dim\operatorname{Inv}_5(V_5)=0,
\qquad
\dim\operatorname{Inv}_6(V_5)=6.
\]

The vanishing of the cubic invariant is expected for odd $\ell$ under $SO(3)$, and full $O(3)$ invariance also eliminates all odd powers. Quartic order is therefore the first nonlinear order at which the shape selection of a real $\ell=5$ field can become nontrivial.

\section{Exact Quartic Reduction}

Let $B_L(A)=P_L(A\otimes A)$ be the projection of the symmetric tensor square of $A\in V_5$ into the angular-momentum channel $L$. Define

\[
K_L(A)=\|B_L(A)\|^2.
\]

The permitted even channels are $L=0,2,4,6,8,10$. Direct Clebsch--Gordan reduction yields the exact relations

\[
K_4=\frac{18}{13}K_0,
\]

\[
K_6=\frac{45}{17}K_0-\frac{143}{170}K_2,
\]

\[
K_8=\frac{880}{247}K_0-\frac{117}{95}K_2,
\]

and

\[
K_{10}=\frac{777}{323}K_0+\frac{693}{646}K_2.
\]

Thus the quartic invariant algebra has only two independent directions, as predicted by the dimension count. For the baseline reduced energy considered in this work, the quartic term therefore becomes

\[
F_4(A)=\alpha K_0(A)+\beta K_2(A),
\]

with

\[
\alpha=\frac{30241998291659}{24945292961264\pi},
\qquad
\beta=\frac{503273130591}{3050785612808\pi}>0.
\]

On the unit sphere $\|A\|=1$, the scalar channel is fixed:

\[
K_0=\frac1{11}.
\]

Therefore

\[
F_4\big|_{S^{10}}=\frac{\alpha}{11}+\beta K_2(A).
\]

Since $K_2(A)=\|B_2(A)\|^2\ge0$, the normalized quartic minimum set is exactly

\[
\mathcal M_4=\{A\in S^{10}:B_2(A)=0\}.
\]

This formula is the central quartic result. It reduces the full quartic selection problem to the vanishing of a five-component $V_2$ constraint.

\section{An Explicit Regular Minimum}

Consider

\[
A_\star=\sqrt{\frac35}\,Y_{50}+\sqrt{\frac25}\,Y_{55}^{c}.
\]

It is normalized and satisfies

\[
B_2(A_\star)=0.
\]

The differential of the quadratic map is

\[
DB_2(A)[\delta A]=2P_2(A\otimes\delta A).
\]

At $A_\star$, an explicit $5\times5$ minor of this Jacobian has determinant

\[
\det M=\frac{224\sqrt{1430}}{2924207}\neq0.
\]

Hence

\[
\operatorname{rank}DB_2(A_\star)=5.
\]

Because the radial direction already belongs to the kernel when restricted to the normalized constraint, the tangent restriction retains rank five. The implicit-function theorem therefore gives

\[
\dim T_{A_\star}\mathcal M_4=10-5=5.
\]

The local quartic minimum is not isolated. It belongs to a five-dimensional minimum manifold.

\section{Rotation Versus Shape}

The action of $SO(3)$ generates three infinitesimal rotational tangent vectors. At $A_\star$ their Gram matrix is

\[
R^\top R=10I_3,
\]

so they are linearly independent and the rotational orbit is locally three-dimensional. Equivariance ensures that all three lie in the quartic kernel. Since the full tangent minimum manifold has dimension five, two dimensions remain after rotations are removed:

\[
T_{A_\star}\mathcal M_4
=
T_{A_\star}(SO(3)\cdot A_\star)
\oplus
N_{\mathrm{shape}},
\]

with

\[
5=3_{\mathrm{rotation}}+2_{\mathrm{shape}}.
\]

The quartic tangent Hessian can be written as

\[
H_4^{\mathrm{tan}}
=
2\beta\,(DB_2|_T)^*(DB_2|_T),
\]

so

\[
\dim\ker H_4^{\mathrm{tan}}=5.
\]

Again, three zero modes are rotations and two are genuine shape directions. A convenient exact basis for the latter is

\[
U=-\frac{\sqrt{14}}{14}Y_{52}^{c}+Y_{53}^{c},
\]

\[
V=\frac{\sqrt{14}}{14}Y_{52}^{s}+Y_{53}^{s},
\]

with

\[
\|U\|^2=\|V\|^2=\frac{15}{14},
\qquad
\langle U,V\rangle=0.
\]

These directions are not gauge artifacts. They change intrinsic shape while remaining quartically degenerate.

\section{What the Quartic Result Means}

The mathematical conclusion is exact and deliberately narrow. Quartic invariant structure cannot distinguish all normalized degree-five shapes at the minimum. Its resolving algebra ends at a five-dimensional minimum manifold, of which two dimensions remain physically distinct after quotienting rotations. This is a concrete demonstration that a mathematically coherent descriptive layer can identify states that a richer descriptive layer may later distinguish.

The result should not be converted into a consciousness analogy. These two shape-flat directions are physical degeneracies of a truncated invariant energy. They are not a hidden experiential sector, a mathematical model of qualia, or evidence that consciousness resides in the kernel of a physical description. Their role in the larger argument is epistemological: they prove that finite descriptive resolution can create genuine equivalence classes without establishing ultimate identity.

The next question is therefore mathematical rather than metaphysical. Does the degree-six invariant algebra contain relational information capable of distinguishing the two quartic shape directions? The answer is yes, and it can be shown constructively.
